\chapter{Analytical Study of Four Canonical Algorithms}
\label{chap:analytical_study}

In accordance with Kenneth H. Rosen's \emph{Discrete Mathematics and Its Applications} (8th Ed., Chapter 3) \cite{rosen2019discrete}, we present a comprehensive analytical investigation of four foundational algorithms. For each algorithm, we provide its formal problem specification, working principle, step-by-step trace on a small sample, pseudocode, authoritative Python implementation, live execution output, and mathematical complexity derivation.

% ==============================================================================
% 6.1 LINEAR SEARCH
% ==============================================================================
\section{Linear Search (Rosen §3.1, Algorithm 2)}
\label{sec:linear_search}

\subsection{Problem Statement \& Objective}
Given a sequence of $n$ elements $S = [a_0, a_1, \dots, a_{n-1}]$ and a target key $x$, locate the index $i$ of the first occurrence of $x$ in $S$. If $x \notin S$, return $-1$.

\subsection{Working Principle}
Linear Search conducts a sequential pass from index $0$ to $n-1$, evaluating $a_i == x$ at each position. It operates without any precondition on the ordering of elements.

\subsection{Step-by-Step Example}
Given sequence $S = [15, 23, 7, 42, 19]$ and target $x = 42$:
\begin{enumerate}
    \item \textbf{Step 0}: Probe $a_0 = 15$. $15 \neq 42$. Continue. (1 comparison)
    \item \textbf{Step 1}: Probe $a_1 = 23$. $23 \neq 42$. Continue. (2 comparisons)
    \item \textbf{Step 2}: Probe $a_2 = 7$. $7 \neq 42$. Continue. (3 comparisons)
    \item \textbf{Step 3}: Probe $a_3 = 42$. $42 == 42$. Match found. Return index $3$.
\end{enumerate}
Total Key Comparisons: $4$.

\subsection{Pseudocode (Rosen §3.1, Alg. 2)}
\begin{lstlisting}[caption={Pseudocode: Linear Search (Rosen §3.1, Algorithm 2)}]
procedure linear_search(x: integer, a1, a2, ..., an: integers)
    i := 1
    while (i <= n and x != a_i)
        i := i + 1
    if i <= n then location := i
    else location := 0
    return location
\end{lstlisting}

\subsection{Python Implementation}
Listing \ref{lst:linear_search} presents the authoritative source code from \texttt{src/algorithms/search.py}.

\begin{lstlisting}[language=Python, caption={Authoritative Implementation of Linear Search}, label={lst:linear_search}]
def linear_search(seq: Sequence[Any], target: Any) -> int:
    for index, item in enumerate(seq):
        if item == target:
            return index
    return -1

def linear_search_instrumented(seq: Sequence[Any], target: Any) -> tuple[int, SearchMetrics]:
    n = len(seq)
    key_comparisons = 0
    equality_comparisons = 0
    loop_evaluations = 0

    for index in range(n):
        loop_evaluations += 1
        equality_comparisons += 1
        key_comparisons += 1
        if seq[index] == target:
            return index, SearchMetrics(
                key_comparisons=key_comparisons,
                equality_comparisons=equality_comparisons,
                ordering_comparisons=0,
                loop_evaluations=loop_evaluations,
            )

    loop_evaluations += 1
    return -1, SearchMetrics(
        key_comparisons=key_comparisons,
        equality_comparisons=equality_comparisons,
        ordering_comparisons=0,
        loop_evaluations=loop_evaluations,
    )
\end{lstlisting}

\subsection{Execution Trace Output}
Executing on sample $S = [15, 23, 7, 42, 19]$ yields:
\begin{verbatim}
Best Case   (target=15) -> Index: 0, Key Comparisons: 1
Middle Case (target=7)  -> Index: 2, Key Comparisons: 3
Worst Case  (target=19) -> Index: 4, Key Comparisons: 5
Missing Case(target=99) -> Index: -1, Key Comparisons: 5
\end{verbatim}

\subsection{Mathematical Complexity Derivation}
\begin{itemize}
    \item \textbf{Best Case}: Target is at index $0$. $T_{\text{best}}(n) = 1 = \Theta(1)$.
    \item \textbf{Worst Case (Successful)}: Target is at index $n - 1$. $T_{\text{worst}}(n) = n = \Theta(n)$.
    \item \textbf{Unsuccessful Search}: Target is absent. Algorithm checks all $n$ elements. $T_{\text{unsuccessful}}(n) = n = \Theta(n)$.
    \item \textbf{Average Case (Successful)}: Under uniform target distribution where $P(x = a_i) = 1/n$:
    $$T_{\text{avg}}(n) = \sum_{i=1}^n i \cdot \frac{1}{n} = \frac{1}{n} \cdot \frac{n(n + 1)}{2} = \frac{n + 1}{2} = \Theta(n)$$
    \item \textbf{Auxiliary Space}: $\Theta(1)$ (in-place).
\end{itemize}

% ==============================================================================
% 6.2 BINARY SEARCH
% ==============================================================================
\section{Binary Search (Rosen §3.1, Algorithm 3)}
\label{sec:binary_search}

\subsection{Problem Statement \& Objective}
Given a pre-sorted sequence $S = [a_0, a_1, \dots, a_{n-1}]$ in non-decreasing order ($a_0 \le a_1 \le \dots \le a_{n-1}$) and a target key $x$, determine an index location $i$ where $a_i == x$, or return $-1$ if $x \notin S$.

\subsection{Working Principle}
Binary Search operates by interval halving. It maintains search boundaries $[\text{low}, \text{high}]$. In each iteration, it inspects midpoint $\text{mid} = \lfloor (\text{low} + \text{high}) / 2 \rfloor$. If $a_{\text{mid}} == x$, it halts. If $a_{\text{mid}} < x$, it restricts search to the upper subsegment $[\text{mid} + 1, \text{high}]$. Otherwise, it restricts search to $[\text{low}, \text{mid} - 1]$.

\subsection{Step-by-Step Example}
Sorted sequence: $S = [10, 20, 30, 40, 50, 60, 70]$, Target: $x = 60$:
\begin{enumerate}
    \item \textbf{Probe 1}: $\text{low}=0, \text{high}=6 \implies \text{mid}=3$. Element $a_3 = 40$.
    \begin{itemize}
        \item Equality: $40 == 60$ (False).
        \item Ordering: $40 < 60$ (True). Narrow to right half: $\text{low} = 4$.
    \end{itemize}
    \item \textbf{Probe 2}: $\text{low}=4, \text{high}=6 \implies \text{mid}=5$. Element $a_5 = 60$.
    \begin{itemize}
        \item Equality: $60 == 60$ (True). Match found at index $5$!
    \end{itemize}
\end{enumerate}
Total Key Comparisons: $1 \text{ (ordering)} + 2 \text{ (equalities)} = 3$.

\subsection{Pseudocode (Rosen §3.1, Alg. 3)}
\begin{lstlisting}[caption={Pseudocode: Binary Search (Rosen §3.1, Algorithm 3)}]
procedure binary_search(x: integer, a1, a2, ..., an: increasing integers)
    i := 1  {left endpoint of search interval}
    j := n  {right endpoint of search interval}
    while i < j
        m := floor((i + j) / 2)
        if x > a_m then i := m + 1
        else j := m
    if x = a_i then location := i
    else location := 0
    return location
\end{lstlisting}

\subsection{Python Implementation}
Listing \ref{lst:binary_search} presents the authoritative source code from \texttt{src/algorithms/search.py}.

\begin{lstlisting}[language=Python, caption={Authoritative Implementation of Binary Search}, label={lst:binary_search}]
def binary_search(seq: Sequence[Any], target: Any, *, check_sorted: bool = False) -> int:
    if check_sorted:
        for i in range(len(seq) - 1):
            if seq[i] > seq[i + 1]:
                raise ValueError("Precondition violated: sequence must be sorted.")
    low = 0
    high = len(seq) - 1
    while low <= high:
        mid = (low + high) // 2
        mid_val = seq[mid]
        if mid_val == target:
            return mid
        elif mid_val < target:
            low = mid + 1
        else:
            high = mid - 1
    return -1
\end{lstlisting}

\subsection{Mathematical Complexity Derivation}
The maximum number of interval halving steps satisfies the recurrence relation:
$$T(n) = T(\lfloor n/2 \rfloor) + \mathcal{O}(1)$$
By the Master Theorem, $T(n) = \Theta(\log n)$.
\begin{itemize}
    \item \textbf{Best Case}: Target is located at initial midpoint. $T_{\text{best}}(n) = 1 = \Theta(1)$ key comparison.
    \item \textbf{Worst Case}: Unsuccessful search or target located at leaf level. Maximum probes is $\lfloor \log_2 n \rfloor + 1$. Each non-terminating probe performs 2 comparisons (1 equality + 1 ordering):
    $$T_{\text{worst}}(n) \le 2(\lfloor \log_2 n \rfloor + 1) = \Theta(\log n)$$
    \item \textbf{Auxiliary Space}: $\Theta(1)$ (iterative).
\end{itemize}

% ==============================================================================
% 6.3 BUBBLE SORT
% ==============================================================================
\section{Bubble Sort (Rosen §3.1, Algorithm 4)}
\label{sec:bubble_sort}

\subsection{Problem Statement \& Objective}
Sort a mutable sequence of $n$ elements in place into non-decreasing order ($a_0 \le a_1 \le \dots \le a_{n-1}$) using adjacent pairwise exchanges.

\subsection{Working Principle}
Bubble Sort performs $n - 1$ fixed outer passes. In pass $i$ (for $i = 0$ to $n - 2$), it compares adjacent elements $a_j$ and $a_{j+1}$ for $j = 0$ to $n - 2 - i$. If $a_j > a_{j+1}$, it swaps them. Per Rosen §3.1 (Algorithm 4), the textbook primary version does \emph{not} employ early termination flags.

\subsection{Step-by-Step Example}
Input array: $A = [4, 2, 5, 1]$ ($n = 4$):
\begin{itemize}
    \item \textbf{Pass 0}: Compare $(4, 2) \to$ Swap $[2, 4, 5, 1]$; Compare $(4, 5) \to$ No swap; Compare $(5, 1) \to$ Swap $[2, 4, 1, 5]$. (3 comps, 2 swaps)
    \item \textbf{Pass 1}: Compare $(2, 4) \to$ No swap; Compare $(4, 1) \to$ Swap $[2, 1, 4, 5]$. (2 comps, 1 swap)
    \item \textbf{Pass 2}: Compare $(2, 1) \to$ Swap $[1, 2, 4, 5]$. (1 comp, 1 swap)
\end{itemize}
Total Comparisons: $3 + 2 + 1 = 6 = \frac{4 \times 3}{2}$. Total Swaps: $4$.

\subsection{Pseudocode (Rosen §3.1, Alg. 4)}
\begin{lstlisting}[caption={Pseudocode: Bubble Sort (Rosen §3.1, Algorithm 4)}]
procedure bubblesort(a1, ..., an: real numbers with n >= 2)
    for i := 1 to n - 1
        for j := 1 to n - i
            if a_j > a_{j+1} then swap a_j and a_{j+1}
\end{lstlisting}

\subsection{Python Implementation}
Listing \ref{lst:bubble_sort} presents the authoritative source code from \texttt{src/algorithms/sorting.py}.

\begin{lstlisting}[language=Python, caption={Authoritative Implementation of Bubble Sort}, label={lst:bubble_sort}]
def bubble_sort(arr: MutableSequence[Any]) -> MutableSequence[Any]:
    n = len(arr)
    for i in range(n - 1):
        for j in range(n - 1 - i):
            if arr[j] > arr[j + 1]:
                arr[j], arr[j + 1] = arr[j + 1], arr[j]
    return arr
\end{lstlisting}

\subsection{Mathematical Complexity Derivation}
\begin{itemize}
    \item \textbf{Comparisons (All Cases)}: For any input of size $n$, the fixed pass structure executes:
    $$\sum_{i=0}^{n-2} (n - 1 - i) = (n - 1) + (n - 2) + \dots + 1 = \frac{n(n - 1)}{2} = \Theta(n^2)$$
    \item \textbf{Best-Case Swaps}: Input already sorted $\implies 0$ swaps.
    \item \textbf{Worst-Case Swaps}: Reverse-sorted distinct array. Every pair is inverted $\implies \frac{n(n-1)}{2}$ swaps.
    \item \textbf{Average-Case Swaps}: Random permutations $\implies \frac{n(n-1)}{4}$ expected swaps.
    \item \textbf{Auxiliary Space}: $\Theta(1)$ (in-place). Stability: Stable.
\end{itemize}

% ==============================================================================
% 6.4 INSERTION SORT
% ==============================================================================
\section{Insertion Sort (Rosen §3.1, Algorithm 5)}
\label{sec:insertion_sort}

\subsection{Problem Statement \& Objective}
Sort a mutable sequence of $n$ elements in place into non-decreasing order by maintaining a sorted prefix and iteratively inserting each subsequent element into its verified position.

\subsection{Working Principle}
For index $j$ from $1$ to $n-1$:
\begin{enumerate}
    \item Store $\text{key} = a_j$.
    \item Backtrack through the sorted prefix $a_{0 \dots j-1}$ using pointer $i = j - 1$.
    \item While $i \ge 0$ and $a_i > \text{key}$, shift element $a_i$ one index to the right ($a_{i+1} = a_i$) and decrement $i$.
    \item Place $\text{key}$ into the resulting slot: $a_{i+1} = \text{key}$.
\end{enumerate}

\subsection{Counting Convention Definition}
\begin{itemize}
    \item \textbf{Right Shift}: Copying an existing element rightward ($a_{i+1} = a_i$).
    \item \textbf{Insertion Assignment}: Placing the key into the vacant slot ($a_{i+1} = \text{key}$). This is \emph{not} counted as a right shift.
\end{itemize}

\subsection{Pseudocode (Rosen §3.1, Alg. 5)}
\begin{lstlisting}[caption={Pseudocode: Insertion Sort (Rosen §3.1, Algorithm 5)}]
procedure insertion_sort(a1, a2, ..., an: real numbers with n >= 2)
    for j := 2 to n
        i := 1
        while aj > ai
            i := i + 1
        m := aj
        for k := 0 to j - i - 1
            aj - k := aj - k - 1
        ai := m
\end{lstlisting}

\subsection{Python Implementation}
Listing \ref{lst:insertion_sort} presents the authoritative source code from \texttt{src/algorithms/sorting.py}.

\begin{lstlisting}[language=Python, caption={Authoritative Implementation of Insertion Sort}, label={lst:insertion_sort}]
def insertion_sort(arr: MutableSequence[Any]) -> MutableSequence[Any]:
    n = len(arr)
    for j in range(1, n):
        key = arr[j]
        i = j - 1
        while i >= 0 and arr[i] > key:
            arr[i + 1] = arr[i]
            i -= 1
        arr[i + 1] = key
    return arr

def insertion_sort_instrumented(arr: MutableSequence[Any]) -> tuple[MutableSequence[Any], SortingMetrics]:
    n = len(arr)
    element_comparisons = 0
    right_shifts = 0

    for j in range(1, n):
        key = arr[j]
        i = j - 1
        while i >= 0:
            element_comparisons += 1
            if arr[i] > key:
                arr[i + 1] = arr[i]
                right_shifts += 1
                i -= 1
            else:
                break
        arr[i + 1] = key

    return arr, SortingMetrics(element_comparisons=element_comparisons, swaps=0, right_shifts=right_shifts)
\end{lstlisting}

\subsection{Mathematical Complexity Derivation}
\begin{itemize}
    \item \textbf{Best Case (Already Sorted)}: In every step $j$, $a_{j-1} > \text{key}$ evaluates to False on the first probe.
    $$T_{\text{best}}(n) = n - 1 = \Theta(n) \text{ comparisons}, \quad 0 \text{ right shifts}$$
    \item \textbf{Worst Case (Reverse Sorted Distinct)}: Key shifts past all $j$ items in prefix.
    $$\text{Comparisons} = \sum_{j=1}^{n-1} j = \frac{n(n - 1)}{2} = \Theta(n^2), \quad \text{Shifts} = \frac{n(n - 1)}{2} = \Theta(n^2)$$
    \item \textbf{Average Case}: $\Theta(n^2)$ comparisons and right shifts ($\approx n^2 / 4$).
    \item \textbf{Auxiliary Space}: $\Theta(1)$ (in-place). Stability: Stable.
\end{itemize}
